\documentclass[11pt]{article}
\usepackage[margin=1in]{geometry}
\usepackage{amsmath,amssymb,amsthm,hyperref}
\hypersetup{colorlinks=true,linkcolor=blue,urlcolor=blue}
\newcommand{\sech}{\operatorname{sech}}
\newcommand{\Res}{\operatorname{Res}}
\newtheorem{theorem}{Theorem}
\title{The first-moment Berndt integral of order three\\MathDB 360892}
\author{Research note}
\date{17 September 2026}
\begin{document}
\maketitle

\section*{Task 1. Prove the exact integral evaluation}
Write $G=\Gamma(1/4)$. We prove Conjecture 4.3 of Rui--Xu--Zhao
\cite{RXZ,RS}, with precisely its stated integrand and constants.
Theorem 4.1 of those sources treats numerators $x^{4m+1}$ for $m>0$;
the present numerator $x$ is outside that range.

\begin{theorem}
\[
 \boxed{\displaystyle
 \int_0^\infty\frac{x\,dx}{(\cos x+\cosh x)^3}
 =-\frac{G^4}{2^7\pi^2}+\frac{G^4}{2^9\pi}
                         +\frac{G^{12}}{2^{13}\pi^7}.}
\]
\end{theorem}

\subsection*{1. A convergent residue formula}
Denote the integral by $I$, and put
\[
 a=\frac{1+i}{2},\qquad b=\frac{1-i}{2},\qquad
 t_n=\pi(n+\tfrac12),\qquad h(t)=\sech^3t.
\]
The integral converges: the integrand is $O(x)$ at zero and
$O(xe^{-3x})$ at infinity. Use
\[
 F(z)=\frac{z}{(\cos z+\cosh z)^3},\qquad
 \cos z+\cosh z=2\cosh(az)\cosh(bz).
\]
On the positively oriented triangle with vertices $0,R,iR$, take
$R=2\pi N$, $N$ a positive integer. For the oblique edge $x+y=R$,
$\operatorname{Im}(az)=\pi N$ and $\operatorname{Re}(bz)=R/2$.
Consequently $|\cosh(az)|\ge1$ and
$|\cosh(bz)|\ge\sinh(R/2)$. Its integral is therefore
$O(R^2e^{-3R/2})$. Also $F(iy)=iF(y)$, so the two axial integrals
sum to $2\int_0^R F(x)\,dx$. The first-quadrant poles are exactly
$z_n=(1+i)t_n$, of order three. Thus
\begin{equation}\label{res-sum}
 I=\pi i\sum_{n\ge0}\Res_{z=z_n}F(z).
\end{equation}

To evaluate a residue, write $z=z_n+v$. Then
\[
 F(z_n+v)=\frac{i(-1)^n}{8a^3v^3}
       (z_n+v)h(t_n+bv)\left(\frac{av}{\sinh(av)}\right)^3,
 \qquad
 \left(\frac{av}{\sinh(av)}\right)^3
             =1-\frac{a^2v^2}{2}+O(v^4).
\]
The coefficient of $v^{-1}$, multiplied by $\pi i$, is
\[
 \pi i\Res_{z=z_n}F(z)
 =\frac{\pi(-1)^n}{8}
       \{t_n h''(t_n)+2h'(t_n)+t_n h(t_n)\}.
\]
These terms decay exponentially, which also justifies the limiting
sum in \eqref{res-sum}. If $s(t)=\sech t$, then $s''=s-2s^3$.
Replacing $h$ by $(s-s'')/2$ gives the particularly useful formula
\begin{equation}\label{reduced}
 I=\frac{\pi}{16}\sum_{n\ge0}(-1)^n
 \{t_n(s-s^{(4)})(t_n)+2(s'-s^{(3)})(t_n)\}.
\end{equation}

\subsection*{2. Interchange the divisor variables}
For nonnegative integers $P,Q$ and $|q|<1$, define
\[
 \Phi_{P,Q}(q)=\sum_{k,v\ge1}k^P v^Q q^{kv},\qquad
 q_*=e^{-\pi/2},\qquad
 T_{P,Q}=\frac{\Phi_{P,Q}(iq_*)-\Phi_{P,Q}(-iq_*)}{2i}.
\]
Let $\chi_4(j)$ be zero for even $j$ and $(-1)^{(j-1)/2}$ for odd
$j$. Since $\chi_4(kv)=\chi_4(k)\chi_4(v)$, we have
\[
 T_{P,Q}=\sum_{k,v\ge1}\chi_4(k)\chi_4(v)k^Pv^Qq_*^{kv}.
\]
The absolutely convergent expansion
$s(t)=2\sum_{v\ge1}\chi_4(v)e^{-vt}$ now yields, for every
nonnegative derivative order $d$,
\begin{equation}\label{twist}
 \sum_{n\ge0}(-1)^nt_n^P s^{(d)}(t_n)
       =2^{1-P}\pi^P(-1)^d T_{P,d}.
\end{equation}
All rearrangements and differentiations here are justified by
absolute convergence; for example $e^{-\pi kv/2}$ dominates every
fixed polynomial weight. Applying \eqref{twist} to \eqref{reduced},
\begin{equation}\label{four-sums}
 I=\frac{\pi^2}{16}(T_{1,0}-T_{1,4})
                   -\frac{\pi}{4}(T_{0,1}-T_{0,3}).
\end{equation}

Use the normalization $q=e^{2\pi i\tau}$ and $D=q\,d/dq$.
For $r\ge1$ put
$L_r(q)=\sum_{k,v\ge1}v^r q^{kv}=\sum_{n\ge1}\sigma_r(n)q^n$.
The standard Eisenstein series satisfy
\[
 E_2=1-24L_1,\qquad E_4=1+240L_3,\qquad E_6=1-504L_5.
\]
Interchanging $k$ and $v$, and differentiating absolutely convergent
series, proves
\begin{align}
 \Phi_{1,0}=\Phi_{0,1}&=\frac{1-E_2}{24},
 &\Phi_{0,3}&=\frac{E_4-1}{240},\label{Phi}\\
 \Phi_{1,4}=DL_3&=\frac{E_2E_4-E_6}{720}.
 &&\label{Phi14}
\end{align}
The last equality is Ramanujan's differential identity
$DE_4=(E_2E_4-E_6)/3$ \cite[Section 5.1]{Z}.
In particular, no series with a negative power of its divisor
variable has to be evaluated.

\subsection*{3. The required modular values}
Put
\[
 u=\frac{G^4}{4\pi^3},\qquad \tau_+=\frac{1+i}{4},\qquad
 \tau_- =\frac{-1+i}{4}.
\]
We will establish
\begin{equation}\label{cm}
 E_2(\tau_+)=\frac{12}{\pi}-3iu,\qquad
 E_4(\tau_+)=-33u^2,\qquad E_6(\tau_+)=189iu^3.
\end{equation}
Here are the details, including the weight-two correction.

Use the standard theta constants
\[
 \theta_2(\tau)=\sum_{n\in\mathbb Z}e^{\pi i\tau(n+1/2)^2},\quad
 \theta_3(\tau)=\sum_{n\in\mathbb Z}e^{\pi i\tau n^2},\quad
 \theta_4(\tau)=\sum_{n\in\mathbb Z}(-1)^ne^{\pi i\tau n^2}.
\]
The classical values are $E_2(i)=3/\pi$ and $E_4(i)=3u^2/4$
\cite[p.~85]{Z}. Theta transformation and Jacobi's identity give
$\theta_2(i)^4=\theta_4(i)^4=\theta_3(i)^4/2$.
Together with the standard identities \cite[pp.~28--29]{Z}
\begin{align*}
 E_4&=\tfrac12(\theta_2^8+\theta_3^8+\theta_4^8),\\
 E_6&=\tfrac12(\theta_2^4+\theta_3^4)
                 (\theta_3^4+\theta_4^4)(\theta_4^4-\theta_2^4),
\end{align*}
and positivity on the imaginary axis, these imply
$\theta_3(i)^4=u$.

The theta duplication identities \cite[Section 20.7]{DLMF} are
\[
 \theta_2(2\tau)^2=\frac{\theta_3(\tau)^2-\theta_4(\tau)^2}{2},
 \quad
 \theta_3(2\tau)^2=\frac{\theta_3(\tau)^2+\theta_4(\tau)^2}{2},
 \quad
 \theta_4(2\tau)^4=\theta_3(\tau)^2\theta_4(\tau)^2.
\]
Thus, at $2i$,
\[
 (\theta_2^4,\theta_3^4,\theta_4^4)
  =u\left(\frac{3-2\sqrt2}{8},\frac{3+2\sqrt2}{8},
                                  \frac1{\sqrt2}\right),
 \qquad
 E_4(2i)=\frac{33u^2}{64},\quad E_6(2i)=\frac{189u^3}{512}.
\]
For $E_2$, the elementary level-two identity is
\begin{equation}\label{level-two}
 2E_2(2\tau)-E_2(\tau)
                    =\frac{\theta_3(\tau)^4+\theta_4(\tau)^4}{2}.
\end{equation}
To see its normalization directly, the coefficient of $q^n$ on the
right is the number $r_4(2n)$ of representations as four squares.
Jacobi's formula \cite[Proposition 11]{Z} gives
\[
 r_4(2n)=8\sum_{d\mid2n,\,4\nmid d}d
       =24\{\sigma_1(n)-2\sigma_1(n/2)\},
\]
where $\sigma_1(n/2)=0$ for odd $n$. These are exactly the
coefficients on the left, and both constant terms are one.
At $\tau=i$, \eqref{level-two} therefore gives
$E_2(2i)=3/(2\pi)+3u/8$.

Finally set $z=2i-2$, so that $\tau_+=-1/z$.
Integer translation leaves the three Eisenstein series unchanged.
Their inversion laws are
\[
 E_2(-1/z)=z^2E_2(z)+\frac{6z}{\pi i},\qquad
 E_4(-1/z)=z^4E_4(z),\qquad E_6(-1/z)=z^6E_6(z).
\]
Substitution of the three values at $2i$ proves \eqref{cm}.

\subsection*{4. Finish the evaluation}
The numbers $iq_*$ and $-iq_*$ correspond to $\tau_+$ and $\tau_-$.
The Fourier coefficients of $E_2,E_4,E_6$ are real, so their values
at $\tau_-$ are the complex conjugates of those in \eqref{cm}.
Consequently \eqref{Phi}--\eqref{Phi14} give
\[
 T_{1,0}=T_{0,1}=\frac{u}{8},\qquad T_{0,3}=0,\qquad
 T_{1,4}=\frac{\operatorname{Im}\{(12/\pi-3iu)(-33u^2)
                                      -189iu^3\}}{720}
          =-\frac{u^3}{8}.
\]
Formula \eqref{four-sums} becomes
\[
 I=\frac{\pi^2}{128}(u+u^3)-\frac{\pi u}{32}.
\]
Replacing $u$ by $G^4/(4\pi^3)$ gives exactly the three terms in
the theorem. This is an identity obtained by convergent analytic
operations and classical modular identities; no numerical fitting
or algebraic-independence assumption is used.\qed

\section*{Task 2. Evaluate all first moments}
Extend the residue and divisor-interchange calculation to
$\int_0^\infty x(\cos x+\cosh x)^{-M}\,dx$ for every positive
integer $M$. Determine the exact period support and resolve the
all-power first-moment conjecture of Xu--Zhao, arXiv:2401.01385v2,
Conjecture 5.8, including the smallest powers.

\section*{Task 3. Resolve the auxiliary series questions}
Evaluate the half-integer hyperbolic sums proposed after Conjecture
4.3 in the Research Square version. Check their printed numerators
and denominator powers against the contour formula, and identify
which individual evaluations are necessary for that alternative proof.

\section*{Task 4. Vary the modular parameter}
For $q_*=e^{-y/2}$ with $y>0$, retain the four twisted sums as functions
of $y$. Classify imaginary quadratic values for which their
Eisenstein evaluations lead to explicit gamma-value formulas, and
find contour integrals that realize these other parameters.

\section*{Task 5. Compute an exact recurrence in the power}
Combine the recurrence for $\sech^M t$ with Ramanujan's differential
equations to obtain a recurrence for the first-moment coefficients.
Determine whether it avoids reconstructing every individual residue
and whether it yields useful sign or growth information.

\begin{thebibliography}{4}
\bibitem{RXZ} H.~Rui, C.~Xu and J.~Zhao,
\emph{Berndt-Type Integrals of Order Three and Series Associated with
Jacobi Elliptic Functions},
\href{https://arxiv.org/abs/2311.15666v2}{arXiv:2311.15666v2}
(28 November 2023), Conjecture 4.3, p.~17.
\bibitem{RS} H.~Rui, C.~Xu and J.~Zhao, same title,
\href{https://doi.org/10.21203/rs.3.rs-3969711/v1}{Research Square version 1}
(22 February 2024), Conjecture 4.3, printed pp.~18--19.
\bibitem{Z} D.~Zagier, \emph{Elliptic Modular Forms and Their Applications},
in \emph{The 1-2-3 of Modular Forms}, Springer (2008), pp.~1--103;
\href{https://people.mpim-bonn.mpg.de/zagier/files/doi/10.1007/978-3-540-74119-0_1/fulltext.pdf}{author-hosted chapter}.
\bibitem{DLMF} NIST Digital Library of Mathematical Functions,
\href{https://dlmf.nist.gov/20.7}{Section 20.7, theta identities}.
\end{thebibliography}
\noindent\textit{Verification and status.}
The accompanying \texttt{verify\_evaluation.py} checks the exact
residue algebra, modular values, and final constants. Independent
numerical diagnostics are reported separately. Both complete source
versions and dated public-claim checks accompany the note. The
authoring agent performed the audit; independent peer review has not
occurred. Tasks 2--5 are proposed directions, not additional solutions.
\end{document}
